\subsection{Proof Sketch}\label{sec:rts-proof}
Our analysis follows a similar template as the analysis of ERADE by \cite{hu2009efficient}. We first introduce the natural extension to $k$ arms of the quantities introduced in the analysis of ERADE:
	\begin{itemize}
		\item $p_{m,k} := \mathbb{E}[X_{m,k} \mid \mathcal{F}_{m-1}]$, the conditional probability of assigning treatment $k$ at step $m$
		\item $M_{n,k} := \sum_{m=1}^{n}(X_{m,k} - p_{m,k})$, a martingale capturing the deviation of the actual assignment from its conditional expectation
		%\item $Q_{n,k} := \sum_{m=1}^{n}X_{m,k}(\xi_{m,k} - \theta_{k})$, a martingale capturing the deviation between the samples from treatment $k$ and their average
		%and denote \(\mathbf{Q_n} := (\mathbf{Q_{n,1}}, \dots, \mathbf{Q_{n,K}})\)
	\end{itemize}
as well as
\[U_{n,k} = \sum_{m=1}^{n-1} \alpha \hat{\rho}_{m,k} + M_{n,k} - n \hat{\rho}_{n,k}\;.\]

Our first contribution is to identify some sufficient conditions under which the conclusions of Theorem~\ref{thm:LLN} and Theorem~\ref{thm:normality-gen-erade} follow.

\begin{lemma} \label{lem:generic-conditions-main} Consider a RAR procedure under which there exists a sequence of random variables $(\ell_{n,k})_{n,k}$ satisfying the following properties:
\begin{enumerate}[label=(\roman*)]
		\item for all $k\in [K]$, \(\left(\ell_{n,k}\right)_{n \geq 1}\) is a non-decreasing sequence and for all \(n \in \mathbb{N}: \,\ell_{n,k} \leq n\)
		\item for all $k\in [K]$, $n \in \mathbb{N}$,
		\begin{equation}\label{eq:prelim-sketch}
			N_{n,k} - n\hat{\rho}_{n,k} \leq 1+N_{\ell_{n,k},k} - \ell_{n,k} \hat{\rho}_{\ell_{n,k},k} + U_{n,k} - U_{\ell_{n,k},k}
		\end{equation}
		\item for all $k\in [K]$, $n \in \mathbb{N}$,
		\begin{equation}\label{eq:condition-sketch}
			N_{\ell_{n,k},k} - \ell_{n,k} \hat{\rho}_{\ell_{n,k},k} = o(\sqrt{n}) \ a.s.
			\quad
		\end{equation}
	\end{enumerate}
Then the conclusions of Theorem~\ref{thm:LLN} and Theorem~\ref{thm:normality-gen-erade} hold.
\end{lemma}

We obtained these conditions by isolating the key arguments that are used in the ERADE analysis for $K=2$ arms, and by showing that they can actually be naturally be extended to $K$ arms. In Lemma~\ref{lem:preliminary} in Appendix~\ref{appnA}, we show that these conditions are satisfied for the random sequence
\[\ell_{n,k} := \max \left\{\, m \leq n ; \, \frac{N_{m,k}}{m} \leq \hat{\rho}_{m,k}  \,\right\}\]
for any $\alpha$RTS design. In words $\ell_{n,k}$ is the last time before \(n\) at which treatment \(k\) is under-sampled: its allocation proportion is less than or equal to the estimated target proportion.

\smallskip

We now briefly sketch the proof of Lemma~\ref{lem:generic-conditions-main}. The detailed proof that Theorem~\ref{thm:LLN} holds can be found in Appendix~\ref{proof:thm1}. The first step to establish the convergence of the empirical proportions is to show that the sequence $\hat{\rho}_n$ has a limit $u \in \Delta_K$ that is non-sparse, i.e. that satisfies $u_k \neq 0$ for all $k$. The argument there relies on both Condition \hyperlink{CA}{\textbf{A}} (to apply the law of large number) and Condition \hyperlink{CB}{\textbf{B}}.

Then, combining Inequalities~\eqref{eq:prelim-sketch} and \eqref{eq:condition-sketch} yields that for all $k$
\begin{equation}\label{step-main}\frac{N_{n,k}}{n} - \hat{\rho}_{n,k} \leq \frac{U_{n,k}- U_{\ell_{n,k},k}}{n} + o\left(\frac{1}{\sqrt{n}}\right).\end{equation}
Starting from the definition of $U_{n,k}$, we can now apply the law of large number to the martingale $M_n$ and use the convergence of $\hat{\rho}_{n}$ to prove that $\lim_{n\rightarrow \infty} \frac{U_{n,k}}{n} = -(1-\alpha)u_k \in (0,1)$. Some algebra (Lemma~\ref{lem:convergence}) further permits to show that the right-hand side of Equation~\eqref{step-main} tends to zero. Using that both $N_n/n$ and $\hat{\rho}_n$ are probability vectors, it follows that $\lim_{n\rightarrow 0} \frac{N_{n,k}}{n} - \hat{\rho}_{n,k} = 0$ and that ${N_{n}}/{n}$ also converges to $u \in (0,1)^K$.

A straightforward consequence is that $\forall k \in [K], N_{n,k} \to \infty$ a.s., hence $\hat{\Theta}_n$ converges to $\Theta$ and by continuity of \(\rho\) we obtain that \(\hat{\rho}_n = \rho(\hat{\Theta}_n) \to \rho(\Theta) = v \in (0,1)^K\), hence for all $k$,
	\[
	\lim_{n \to \infty} \frac{N_{n,k}}{n} = \lim_{n \to \infty} \hat{\rho}_{n,k} = v_k\;.
	\]
	The convergence rates follows from the Law of Iterated Logarithm.

	\smallskip

	The detailed proof that Theorem~\ref{thm:normality-gen-erade} holds can be found in Appendix~\ref{proof:thm2}. To get the more precise rates in statement (i), we need a more careful control on the increment of the auxiliary process $U_{n,k}$. Namely, we prove through a series of (rather technical) lemmas that
	\[
	U_{n,k}-U_{\ell_{n,k},k}\leq o_p(\sqrt n),
	\qquad
	U_{n,k}-U_{\ell_{n,k},k}\leq O(\sqrt{n\log\log n}) \quad a.s.
	\]
	which yields the following using the identities $\sum_{k=1}^K N_{n,k}=n$ and $\sum_{k=1}^K \hat\rho_{n,k}=1$:
	\[
	N_{n,k}-n\hat\rho_{n,k}=o_p(\sqrt n)\;.
	\]
    This proves the statements \eqref{eq:thm-res-1}--\eqref{eq:thm-res-2}. 	Combined with the rate \(n(\hat\rho_n-v)=O(\sqrt{n\log\log n})\)
	from Theorem \ref{thm:LLN} yields \eqref{eq:thm-res-3}.
	
	To establish the asymptotic normality in (ii), we prove the following lemma that gives the joint CLT of the components of the estimator \(\hat{\Theta}_n\) scaled with \(\sqrt{N_{n,k}}\). The proof of Lemma \ref{lem:componentwise} relies on a martingale-CLT argument and is given in Appendix~\ref{proof:thm2}.
	
		\begin{restatable}{lemma}{jointcltlemma}[Joint CLT with Componentwise Scaling]\label{lem:componentwise} 
			Consider a RAR procedure under which $N_{n,k}\to\infty$ a.s. for all $k \in [K]$, then we have
			\[
			D_n(\hat\Theta_n-\Theta)\xrightarrow{\mathcal{D}}
			\mathcal N\!\big(0,\mathrm{diag}\left(V_1,\dots,V_K\right)\big)
			\]
			where $D_n:=\mathrm{diag}(\sqrt{N_{n,1}},\dots,\sqrt{N_{n,K}})\;.$
		\end{restatable}
		
		Using the fact that \(\sqrt n D_n^{-1}\) converges to \(\mathrm{diag}\left(\frac{1}{\sqrt{v_1}},\ldots,\frac{1}{\sqrt{v_K}}\right)\), the CLT result of Lemma \ref{lem:componentwise} combined with Slutsky's lemma yields the first asymptotic normality result:
		\begin{align*}
			\sqrt n(\hat\Theta_n-\Theta) &= \sqrt n D_n^{-1} D_n(\hat\Theta_n-\Theta)\\
			&\xrightarrow{\mathcal{D}}\mathrm{diag}\left(\frac{1}{\sqrt{v_1}},\ldots,\frac{1}{\sqrt{v_K}}\right) \mathcal N\!\big(0,\mathrm{diag}\left(V_1,\dots,V_K\right)\big)\\ 
			&= \mathcal N(0,V).
		\end{align*}

	
	A first-order Taylor expansion of the target function around $\Theta$ gives
	\[
	\hat\rho_n-v
	= G(\hat\Theta_n-\Theta) + O_p(n^{-1/2}),
	\]
	and hence
	\[
	\sqrt n(\hat\rho_n-v)\xrightarrow{\mathcal{D}}\mathcal N(0,GVG^\top).
	\]
	Finally, the bound \eqref{eq:thm-res-1} implies
	\[
	\sqrt n\!\left(\frac{N_n}{n}-\hat\rho_n\right)=o_p(1),
	\]
	so Slutsky’s lemma yields the CLT for $\sqrt n(N_n/n-v)$ as well as the joint convergence stated in \eqref{eq:thm-res-4} and \eqref{eq:thm-res-5}.
	
	\begin{remark}
			We emphasize that one component of the proof, Lemma \ref{lem:componentwise}, only requires that each treatment is sampled
			infinitely often. The normalization in this CLT result is performed componentwise through
			\(D_n\), allowing each estimator \(\hat\theta_{n,k}\) to be scaled according
			to its own sample size \(N_{n,k}\). As a result, this Lemma provides a general asymptotic normality result for the treatment effect estimators that
			does not rely on any assumption on the limiting allocation proportions. This result will be particularly relevant to tackle sparse target allocations, as we do in the next section.
	\end{remark}
	
